% Vietnamese translation for OI Public Library
% Source-PDF: IOI中国国家候选队论文/2002/张家琳.pdf
\documentclass[11pt]{article}
\usepackage{oipl-vn}

\title{Nhân đa thức nhanh}
\author{Zhang Jialin}
\date{}

\newcommand{\eps}{\varepsilon}
\newcommand{\ii}{\mathrm{i}}
\newcommand{\poly}{\mathcal{P}}

\begin{document}
\maketitle

\origtitle{Fast polynomial multiplication}
\sourcepath{IOI China candidate team papers / 2002 / Zhang Jialin.pdf}

\tableofcontents

\section{Dẫn nhập}

Đa thức là một trong những công cụ toán học cơ bản nhất. Vì có dạng đơn giản
và dễ thực hiện các phép tính trên máy tính, đa thức ngày càng được dùng rộng
rãi. Nó không chỉ giữ vai trò quan trọng trong các phần mềm toán học như
Maple, mà còn có nhiều ứng dụng trong kỹ thuật, thông tin và nhiều lĩnh vực
khác.

Sau đây là một vài ví dụ về xấp xỉ bằng đa thức:
\[
  \ln(x+1)=x-\frac{x^2}{2}+\frac{x^3}{3}-\frac{x^4}{4}
    +\cdots+(-1)^{n-1}\frac{x^n}{n}+\cdots,
\]
\[
  \sin x=x-\frac{x^3}{3!}+\frac{x^5}{5!}-\frac{x^7}{7!}
    +\cdots+(-1)^n\frac{x^{2n+1}}{(2n+1)!}+\cdots,
\]
\[
  \cos x=1-\frac{x^2}{2!}+\frac{x^4}{4!}-\frac{x^6}{6!}
    +\cdots+(-1)^n\frac{x^{2n}}{(2n)!}+\cdots.
\]

\section{Các phép toán cơ bản trên đa thức}

Các phép toán cơ bản gồm phép cộng, phép tính giá trị tại một điểm và phép
nhân. Với
\[
  f(x)=a_0+a_1x+a_2x^2+\cdots+a_nx^n,
\]
ta có thể tính giá trị theo lược đồ Horner:
\[
  f(x)=a_0+x\bigl(a_1+x(a_2+\cdots+x(a_{n-1}+xa_n)\cdots)\bigr).
\]

Phần còn lại tập trung vào phép nhân đa thức.

\section{Phép nhân đa thức thông thường}

Nhân đa thức là một bài toán rất thường gặp. Với thuật toán thông thường,
phép nhân hai đa thức bậc \(n\) cần thời gian \(O(n^2)\). Thuật toán ấy có
dạng quen thuộc như phép nhân đặt cột:
\[
\begin{array}{rrrrr}
&& a_nx^n+a_{n-1}x^{n-1}+\cdots+a_1x+a_0\\
&\times& b_nx^n+b_{n-1}x^{n-1}+\cdots+b_1x+b_0\\ \hline
&& a_nb_0x^n+a_{n-1}b_0x^{n-1}+\cdots+a_1b_0x+a_0b_0\\
&+& a_nb_1x^{n+1}+a_{n-1}b_1x^n+\cdots+a_1b_1x^2+a_0b_1x\\
&& \cdots\\
&+& a_nb_nx^{2n}+\cdots+a_1b_nx^{n+1}+a_0b_nx^n.
\end{array}
\]
Vì thế tích \(r(x)=f(x)g(x)\) có các hệ số
\[
  r(x)=a_nb_nx^{2n}
  +\cdots+
  \sum_{k=1}^{n} a_{n+1-k}b_k\,x^{n+1}
  +\sum_{k=0}^{n} a_{n-k}b_k\,x^n
  +\cdots+(a_1b_0+a_0b_1)x+a_0b_0.
\]

Nhìn vào quá trình trên, có vẻ thuật toán không làm việc thừa: ta phải xét
mọi cặp \(a_i b_j\), vì mỗi tích đều ảnh hưởng đến kết quả cuối cùng. Nếu
chỉ thay đổi thứ tự tính toán, ta không thể giảm bậc độ phức tạp thời gian;
cùng lắm chỉ giảm được hằng số.

Do đó, muốn giảm độ phức tạp một cách căn bản, ta phải nhìn đa thức từ một
góc khác.

\section{Biểu diễn điểm--giá trị của đa thức}

Bên cạnh biểu diễn bằng hệ số, ta có thể biểu diễn đa thức bằng
\term{điểm--giá trị}. Cụ thể, dùng
\[
  (x_0,y_0),(x_1,y_1),(x_2,y_2),\ldots,(x_{n-1},y_{n-1}),
\]
trong đó \(f(x_i)=y_i\) và các \(x_i\) đôi một khác nhau, để biểu diễn một đa
thức bậc không quá \(n-1\).

Với một đa thức đã cho, cách đơn giản nhất để sinh \(n\) cặp điểm--giá trị
là chọn tùy ý \(n\) giá trị \(x_i\) phân biệt, rồi lần lượt tính giá trị đa
thức tại các điểm ấy.

Ngược lại, \(n\) cặp điểm--giá trị cũng xác định duy nhất một đa thức bậc
không quá \(n-1\). Quá trình khôi phục đa thức từ các điểm gọi là
\term{nội suy}.

\paragraph{Bổ đề 1: tính duy nhất của nội suy đa thức.}
Với tập bất kỳ gồm \(n\) cặp điểm--giá trị
\[
  \{(x_0,y_0),(x_1,y_1),\ldots,(x_{n-1},y_{n-1})\},
\]
trong đó \(x_0,x_1,\ldots,x_{n-1}\) đôi một khác nhau, tồn tại duy nhất một
đa thức \(A(x)\) bậc không quá \(n-1\) sao cho
\[
  y_k=A(x_k),\qquad k=0,1,\ldots,n-1.
\]

\paragraph{Chứng minh tồn tại.}
Đặt
\[
  A(x)=a_0+a_1x+a_2x^2+\cdots+a_{n-1}x^{n-1}.
\]
Điều kiện \(A(x_k)=y_k\) tương đương hệ phương trình tuyến tính
\[
\begin{bmatrix}
1 & x_0 & x_0^2 & \cdots & x_0^{n-1}\\
1 & x_1 & x_1^2 & \cdots & x_1^{n-1}\\
\vdots & \vdots & \vdots & \ddots & \vdots\\
1 & x_{n-1} & x_{n-1}^2 & \cdots & x_{n-1}^{n-1}
\end{bmatrix}
\begin{bmatrix}
a_0\\ a_1\\ \vdots\\ a_{n-1}
\end{bmatrix}
=
\begin{bmatrix}
y_0\\ y_1\\ \vdots\\ y_{n-1}
\end{bmatrix}.
\]
Ký hiệu hệ này là \(XA=Y\). Ma trận \(X\) là ma trận Vandermonde. Bằng các
phép biến đổi định thức, ta có
\[
  \det X=\prod_{j<k}(x_k-x_j)\ne0.
\]
Vì vậy \(X\) khả nghịch và \(A=X^{-1}Y\) tồn tại.

\paragraph{Chứng minh duy nhất.}
Giả sử hai đa thức \(f(x)\) và \(g(x)\), đều có bậc không quá \(n-1\), cùng
thỏa các cặp điểm--giá trị trên. Khi đó
\[
  r(x)=f(x)-g(x)
\]
có \(n\) nghiệm \(x_0,x_1,\ldots,x_{n-1}\), trong khi \(r(x)\) có bậc không
quá \(n-1\). Suy ra \(r(x)\equiv0\), tức \(f(x)=g(x)\).

\section{Nhân đa thức trong biểu diễn điểm--giá trị}

Nếu \(r(x)=f(x)g(x)\), thì tại từng điểm \(x_i\),
\[
  r(x_i)=f(x_i)g(x_i).
\]
Nói cách khác, nếu ta có ba biểu diễn điểm--giá trị dùng cùng các hoành độ,
\[
  (x_i,r_i),\qquad (x_i,f_i),\qquad (x_i,g_i),
\]
thì
\[
  r_i=f_i g_i.
\]

Vì vậy, nếu dùng đúng cách biểu diễn điểm--giá trị, phép nhân đa thức chỉ
cần thời gian tuyến tính.

Tuy nhiên, nếu \(f\) và \(g\) đều có bậc không quá \(n-1\), thì \(r=fg\) có
bậc không quá \(2n-2\). Để xác định duy nhất \(r(x)\), ta cần \(2n-1\) cặp
điểm--giá trị, trong khi ban đầu chỉ có \(n\) cặp. Ta giải quyết bằng cách
mở rộng số cặp điểm--giá trị của \(f(x)\) và \(g(x)\) ngay từ đầu lên
\(2n-1\).

\section{Cải thiện phép nhân hệ số bằng biểu diễn điểm--giá trị}

Ta muốn tận dụng việc nhân trong biểu diễn điểm--giá trị chỉ cần thời gian
tuyến tính để tăng tốc phép nhân đa thức vốn được cho bằng hệ số. Cần làm
ba bước:
\begin{enumerate}
  \item chuyển từ biểu diễn hệ số sang biểu diễn điểm--giá trị, gọi là quá
  trình tính giá trị;
  \item nhân từng điểm trong biểu diễn điểm--giá trị;
  \item chuyển từ biểu diễn điểm--giá trị về biểu diễn hệ số, gọi là nội suy.
\end{enumerate}
Bước thứ hai chỉ cần thời gian tuyến tính. Mấu chốt là bước thứ nhất và bước
thứ ba.

\section{Từ hệ số sang điểm--giá trị}

Ta cần tính
\[
  x_i\mapsto f(x_i),\qquad i=0,1,\ldots,n-1.
\]
Nếu dùng Horner trực tiếp thì
\[
  f(x_i)=((\cdots(a_{n-1}x_i+a_{n-2})x_i+a_{n-3})\cdots+a_1)x_i+a_0.
\]
Điểm quan trọng là các \(x_0,x_1,\ldots,x_{n-1}\) do ta tự chọn. Bằng cách
chọn phù hợp, ta có thể đưa quá trình chuyển đổi xuống \(O(n\log n)\).

Ta chọn các căn bậc \(n\) của \(1\):
\[
  \eps_n^0,\eps_n^1,\ldots,\eps_n^{n-1},
\]
trong đó
\[
  \eps_n^k=e^{2\pi\ii k/n}
  =\cos\frac{2\pi k}{n}+\ii\sin\frac{2\pi k}{n}.
\]

\paragraph{Bổ đề 2: định lý chia đôi.}
Với mọi số nguyên \(n\ge0\) và \(k\ge0\), ta có
\[
  \left(\eps_{2n}^{2k}\right)=\eps_n^k.
\]

\paragraph{Chứng minh.}
\[
  \eps_{2n}^{2k}
  =\cos\frac{2\pi\cdot2k}{2n}
   +\ii\sin\frac{2\pi\cdot2k}{2n}
  =\cos\frac{2\pi k}{n}
   +\ii\sin\frac{2\pi k}{n}
  =\eps_n^k.
\]

Bài toán cần giải là tính
\[
  f(\eps_n^k)=a_0+a_1\eps_n^k+\cdots+a_{n-1}(\eps_n^k)^{n-1},
  \qquad k=0,1,\ldots,n-1.
\]
Giả sử \(n\) chẵn; nếu không, ta thêm các hệ số bậc cao bằng \(0\) để đưa
\(n\) về số chẵn. Đặt
\[
  f^{[0]}(x)=a_0+a_2x+\cdots+a_{n-2}x^{n/2-1},
\]
\[
  f^{[1]}(x)=a_1+a_3x+\cdots+a_{n-1}x^{n/2-1}.
\]
Đa thức \(f^{[0]}\) chứa mọi hệ số chỉ số chẵn của \(f\), còn \(f^{[1]}\)
chứa mọi hệ số chỉ số lẻ. Khi đó
\[
  f(x)=f^{[0]}(x^2)+x f^{[1]}(x^2).
\]
Vì vậy, để tính \(f\) tại \(\eps_n^0,\eps_n^1,\ldots,\eps_n^{n-1}\), ta cần
tính \(f^{[0]}\) và \(f^{[1]}\) tại các điểm
\[
  (\eps_n^0)^2,(\eps_n^1)^2,\ldots,(\eps_n^{n-1})^2.
\]

Theo định lý chia đôi,
\[
\begin{aligned}
  (\eps_n^0)^2,(\eps_n^1)^2,\ldots,(\eps_n^{n-1})^2
  &=
  \eps_{n/2}^0,\eps_{n/2}^1,\ldots,\eps_{n/2}^{n-1}\\
  &=
  \eps_{n/2}^0,\eps_{n/2}^1,\ldots,\eps_{n/2}^{n/2-1},
  \eps_{n/2}^0,\ldots,\eps_{n/2}^{n/2-1}.
\end{aligned}
\]
Nghĩa là các bình phương trên không phải \(n\) số khác nhau; chúng chỉ gồm
các căn bậc \(n/2\) của \(1\), mỗi căn xuất hiện đúng hai lần. Do đó bài
toán con có cùng dạng với bài toán gốc nhưng kích thước giảm một nửa. Điều
này dẫn đến cách giải chia để trị, chính là ý tưởng của FFT.

\paragraph{Thuật toán đệ quy.}
Trong đoạn giả mã sau, \(a\) là dãy hệ số, \(y\) là dãy giá trị tại các căn
đơn vị tương ứng.
\begin{verbatim}
Function transform(a: atype): ytype;
  if n = 1 then
    return a0

  if odd(n) then
    begin
      inc(n);
      a[n-1] := 0;
    end;

  a[0] := (a0, a2, ..., a[n-2]);
  a[1] := (a1, a3, ..., a[n-1]);

  y[0] := transform(a[0]);
  y[1] := transform(a[1]);

  for k := 0 to n-1 do
    y[k] := y[0][k mod (n/2)]
            + eps_n^k * y[1][k mod (n/2)];

  return y;
\end{verbatim}
Trong công thức ở vòng lặp,
\[
  y(\eps_n^k)
  =y^{[0]}(\eps_{n/2}^k)+\eps_n^k y^{[1]}(\eps_{n/2}^k),
\]
với chỉ số ở vế phải hiểu theo chu kỳ \(n/2\). Có thể chứng minh cách tính
này có độ phức tạp \(O(n\log n)\).

\section{Từ điểm--giá trị về hệ số}

Nội suy là phép ngược của quá trình tính giá trị. Bài toán này có vẻ phức
tạp hơn, nhưng thực ra có thể chuyển về bài toán trước.

Quá trình tính giá trị có thể viết thành phương trình ma trận
\[
  Y=V_nA,
\]
trong đó
\[
Y=
\begin{bmatrix}
y_0\\ y_1\\ y_2\\ y_3\\ \vdots\\ y_{n-1}
\end{bmatrix},
\qquad
A=
\begin{bmatrix}
a_0\\ a_1\\ a_2\\ a_3\\ \vdots\\ a_{n-1}
\end{bmatrix},
\]
và
\[
V_n=
\begin{bmatrix}
1 & 1 & 1 & 1 & \cdots & 1\\
1 & \eps_n & \eps_n^2 & \eps_n^3 & \cdots & \eps_n^{n-1}\\
1 & \eps_n^2 & \eps_n^4 & \eps_n^6 & \cdots & \eps_n^{2(n-1)}\\
1 & \eps_n^3 & \eps_n^6 & \eps_n^9 & \cdots & \eps_n^{3(n-1)}\\
\vdots & \vdots & \vdots & \vdots & \ddots & \vdots\\
1 & \eps_n^{n-1} & \eps_n^{2(n-1)} & \eps_n^{3(n-1)}
  & \cdots & \eps_n^{(n-1)(n-1)}
\end{bmatrix}.
\]
Nếu \(V_n\) khả nghịch, nội suy là
\[
  A=V_n^{-1}Y.
\]

\paragraph{Bổ đề 3.}
Với ma trận \(V_n\) ở trên,
\[
V_n^{-1}=\frac1n
\begin{bmatrix}
1 & 1 & 1 & 1 & \cdots & 1\\
1 & \eps_n^{-1} & \eps_n^{-2} & \eps_n^{-3}
  & \cdots & \eps_n^{-(n-1)}\\
1 & \eps_n^{-2} & \eps_n^{-4} & \eps_n^{-6}
  & \cdots & \eps_n^{-2(n-1)}\\
1 & \eps_n^{-3} & \eps_n^{-6} & \eps_n^{-9}
  & \cdots & \eps_n^{-3(n-1)}\\
\vdots & \vdots & \vdots & \vdots & \ddots & \vdots\\
1 & \eps_n^{-(n-1)} & \eps_n^{-2(n-1)}
  & \eps_n^{-3(n-1)} & \cdots & \eps_n^{-(n-1)(n-1)}
\end{bmatrix}.
\]

Từ bổ đề này, ta giải được nội suy bằng cùng một thủ tục như quá trình tính
giá trị, chỉ thay các căn \(\eps_n^k\) bằng \(\eps_n^{-k}\), rồi chia mọi hệ
số thu được cho \(n\):
\[
  A=\frac1n V_n(\eps_n^{-1})Y.
\]

\paragraph{Thuật toán nội suy.}
\begin{verbatim}
Function inverse_transform(y: ytype): atype;
  if n = 1 then
    return y0

  if odd(n) then
    begin
      inc(n);
      y[n-1] := 0;
    end;

  y[0] := (y0, y2, ..., y[n-2]);
  y[1] := (y1, y3, ..., y[n-1]);

  a[0] := inverse_transform(y[0]);
  a[1] := inverse_transform(y[1]);

  for k := 0 to n-1 do
    a[k] := a[0][k mod (n/2)]
            + eps_n^(-k) * a[1][k mod (n/2)];

  after the recursion, divide every a[k] by n;
  return a;
\end{verbatim}

\section{Luồng thuật toán nhân đa thức}

Bài toán: \(f(x)\) và \(g(x)\) là hai đa thức bậc không quá \(n-1\), đã biết
biểu diễn hệ số của chúng. Hãy tính biểu diễn hệ số của
\[
  r(x)=f(x)g(x).
\]

Luồng thuật toán:
\begin{enumerate}
  \item \term{Tiền xử lý.} Thêm \(n-1\) hệ số bậc cao bằng \(0\), để bậc có
  thể đạt đến \(2n-2\). Mục đích là có đủ số cặp điểm--giá trị để xác định
  duy nhất \(r(x)\).
  \item \term{Tính giá trị.} Dùng chia để trị qua hàm \texttt{transform} để
  tính giá trị của \(f(x)\) và \(g(x)\) tại các căn bậc \(2n-1\) của \(1\).
  Trong cài đặt FFT thường đệm tiếp đến một lũy thừa của \(2\) không nhỏ hơn
  \(2n-1\).
  \item \term{Nhân từng điểm.} Với mỗi điểm, nhân giá trị của \(f\) và \(g\)
  để thu được giá trị của \(r\) tại điểm đó.
  \item \term{Nội suy.} Hoán đổi vai trò của \(a\) và \(y\), dùng lại
  \texttt{transform} với căn nghịch đảo để tìm biểu diễn hệ số của \(r(x)\),
  rồi chia kết quả cho số điểm.
\end{enumerate}

Bước 1 và bước 3 đều chạy trong \(O(n)\). Bước 2 và bước 4 đều chạy trong
\(O(n\log n)\).

\section{Cải tiến thuật toán}

Trong hàm \texttt{transform}, ta đang dùng đệ quy. Với \(n=8\), quá trình
gọi đệ quy tách các hệ số như sau:
\[
\begin{array}{c}
(a_0,a_1,a_2,a_3,a_4,a_5,a_6,a_7)\\[0.6em]
\swarrow\hspace{8em}\searrow\\[-0.2em]
(a_0,a_2,a_4,a_6)\hspace{4em}(a_1,a_3,a_5,a_7)\\[0.6em]
\swarrow\hspace{2em}\searrow\hspace{5em}\swarrow\hspace{2em}\searrow\\[-0.2em]
(a_0,a_4)\quad(a_2,a_6)\quad(a_1,a_5)\quad(a_3,a_7)\\[0.6em]
\downarrow\hspace{3em}\downarrow\hspace{3em}\downarrow\hspace{3em}\downarrow\\[-0.2em]
(a_0)\quad(a_4)\quad(a_2)\quad(a_6)\quad(a_1)\quad(a_5)\quad(a_3)\quad(a_7).
\end{array}
\]

Phương pháp đệ quy cần nhiều bộ nhớ, vì mỗi lần gọi lại cần truyền một mảng
hệ số mới vào quá trình con. Từ sơ đồ trên, ta thấy có thể làm ngược lại:
thực hiện từ dưới lên bằng phép lặp. Thứ tự lá là
\[
  (a_0),(a_4),(a_2),(a_6),(a_1),(a_5),(a_3),(a_7),
\]
rồi lần lượt ghép thành
\[
  (a_0,a_4),\ (a_2,a_6),\ (a_1,a_5),\ (a_3,a_7),
\]
tiếp đến
\[
  (a_0,a_2,a_4,a_6),\qquad (a_1,a_3,a_5,a_7),
\]
và cuối cùng trở lại
\[
  (a_0,a_1,a_2,a_3,a_4,a_5,a_6,a_7).
\]

\paragraph{Thuật toán lặp.}
\begin{verbatim}
Function transform_better(a: atype): ytype;
  preprocess:
    add high-degree zero coefficients until the length is a power of 2

  initialize:
    y := a

  iterate:
    for k := 1 to lg n do
      merge suitable adjacent blocks and put the result
      into the corresponding positions of y

  return y
\end{verbatim}

Với \(n=8\), các tầng lặp có thể mô tả như sau:
\[
\begin{array}{c}
(a_0)\ (a_4)\ (a_2)\ (a_6)\ (a_1)\ (a_5)\ (a_3)\ (a_7)\\
\eps_1^0\quad \eps_1^0\quad \eps_1^0\quad \eps_1^0\quad
\eps_1^0\quad \eps_1^0\quad \eps_1^0\quad \eps_1^0\\[0.8em]
(a_0,a_4)\qquad (a_2,a_6)\qquad (a_1,a_5)\qquad (a_3,a_7)\\
\eps_2^0,\eps_2^1\qquad \eps_2^0,\eps_2^1\qquad
\eps_2^0,\eps_2^1\qquad \eps_2^0,\eps_2^1\\[0.8em]
(a_0,a_2,a_4,a_6)\qquad\qquad (a_1,a_3,a_5,a_7)\\
\eps_4^0,\eps_4^1,\eps_4^2,\eps_4^3\qquad
\eps_4^0,\eps_4^1,\eps_4^2,\eps_4^3\\[0.8em]
(a_0,a_1,a_2,a_3,a_4,a_5,a_6,a_7)\\
\eps_8^0,\eps_8^1,\eps_8^2,\eps_8^3,\eps_8^4,\eps_8^5,\eps_8^6,\eps_8^7.
\end{array}
\]

\section{So sánh thực nghiệm}

Cuối cùng, ta so sánh phương pháp được giới thiệu trong bài với phép nhân
đa thức thông thường.

Môi trường thử nghiệm: PIII 500, 128M RAM, FreePascal 1.0.4.

\begin{center}
\small
\setlength{\tabcolsep}{4pt}
\begin{tabular}{c|rrrrr}
Bậc đa thức & \(n=100\) & \(n=1000\) & \(n=10000\) & \(n=20000\) & \(n=30000\)\\ \hline
Thường (s) & 0.02 & 0.07 & 5.02 & 18.06 & \(>30\)\\
Điểm--giá trị (s) & 0.05 & 0.50 & 2.46 & 2.88 & 3.36\\
Chính xác \texttt{real} & 10--11 chữ số & 10--11 chữ số & 9--10 chữ số & 8--9 chữ số & 8--9 chữ số
\end{tabular}
\end{center}

Bảng cho thấy khi bậc còn nhỏ, chi phí hằng số của phương pháp điểm--giá trị
có thể lớn hơn. Khi bậc tăng, độ phức tạp \(O(n\log n)\) trở nên vượt trội
so với \(O(n^2)\), dù cần chú ý sai số số thực trong tính toán với căn đơn
vị phức.

\end{document}
